% ============================================================
% Capítulo 8: Conclusiones
% ============================================================
\chapter{Conclusiones}
\label{ch:conclusiones}

Esta investigación examinó las trayectorias académicas de estudiantes de Matemáticas y Física
de la Facultad de Ciencias con el objetivo de caracterizar, explicar y predecir el rezago
académico mediante una arquitectura analítica de tres enfoques complementarios. Los resultados
obtenidos convergen en una narrativa coherente y cuantitativamente robusta que se sintetiza
a continuación.

% ──────────────────────────────────────────────────────────────────────────────
\section{Conclusión Principal: La Paradoja del Perfeccionismo}
% ──────────────────────────────────────────────────────────────────────────────

El hallazgo más significativo y contraintuitivo de esta investigación es la \textbf{paradoja
del perfeccionismo}: los estudiantes con promedio académico alto pero baja regularidad
(\emph{Perfeccionistas}) tienen trayectorias de egreso significativamente más largas que
aquellos con promedio modesto pero alta regularidad continua (\emph{Pragmáticos}).

Esta paradoja fue confirmada de manera convergente por tres metodologías independientes:

\begin{enumerate}
    \item \textbf{Análisis de Supervivencia (Kaplan-Meier y Cox):} Los estudiantes Pragmáticos
          (regularidad alta, promedio bajo al S4) tienen una mediana de egreso de
          $10$--$12$ semestres, mientras que los Perfeccionistas (regularidad baja, promedio alto)
          no alcanzan la mediana de egreso dentro del horizonte de observación de 20 semestres.
          El modelo de Cox cuantificó este efecto: el \emph{Hazard Ratio} de la regularidad
          fue $\text{HR} = 12.70$ ($p < 0.005$), mientras que el promedio natural no fue
          estadísticamente significativo ($p = 0.08$).

    \item \textbf{Modelos Predictivos de Clasificación (Regresión Logística):} La regularidad del
          cuarto semestre (\texttt{s4\_reg}) tuvo el mayor coeficiente beta positivo en el Modelo
          de Alerta Temprana, y la del octavo (\texttt{s8\_reg}) en el Modelo de Trayectoria
          Completa. Ningún predictor de promedio individual figuró entre las covariables de mayor impacto.

    \item \textbf{Regresión Regularizada (Lasso/Ridge):} El único coeficiente con $|t| > 10$
          en el Enfoque Truncado fue \texttt{s8\_reg} ($\beta = -2.445$, $t = -13.64$,
          $p < 10^{-40}$). Los predictores de promedio (natural y real) no resultaron
          significativos de manera individual bajo penalización regularizada.
\end{enumerate}

\noindent
La interpretación es clara: la estrategia de \emph{sacrificar} materias para mantener un
promedio alto —sin priorizar la continuidad en la inscripción— genera una \textbf{deuda de
regularidad} que se acumula y eventualmente produce rezago. La inercia de inscripción
continua, y no la excelencia en la calificación, es el motor estadístico del egreso oportuno
en las licenciaturas de Ciencias exactas.

% ──────────────────────────────────────────────────────────────────────────────
\section{Conclusiones por Capítulo}
% ──────────────────────────────────────────────────────────────────────────────

\subsection{Metodología y Limpieza de Datos (Capítulo 4)}

La metodología de \textbf{Corrección Dinámica a Tiempo Activo} demostró ser superior a la
corrección fija promedio, con un Error Absoluto Medio de la corrección fija de $0.62$
semestres en Matemáticas (máximo: $5.57$ semestres). La reincorporación de estudiantes con
Baja Temporal —excluyendo solo aquellos con más de 4 semestres inactivos— generó un impacto
cuantitativo masivo: la población analítica de Matemáticas creció un $\mathbf{57.4\,\%}$
(de 2{,}687 a 4{,}229 alumnos) y la de Física un $\mathbf{69.3\,\%}$ (de 3{,}013 a
5{,}100 alumnos). Esta decisión metodológica fortaleció sustancialmente la validez estadística
de todas las inferencias posteriores.

\subsection{Análisis Exploratorio (Capítulo 5)}

El PCA confirmó que la alta dimensionalidad del Vector Excalibur (32 variables) puede reducirse
a cinco componentes que explican el $88.5\,\%$ (Matemáticas) y $90.9\,\%$ (Física) de la
varianza total. El \textbf{PC1} actúa como un índice de rendimiento académico global; el
\textbf{PC2} como un índice de trayectoria temporal (mejora o deterioro). La proyección en el
espacio PC1--PC2 revela una gradación continua y estadísticamente significativa que separa a
los estudiantes sin rezago de los que lo experimentan desde etapas tempranas de la carrera.

\noindent
La variable individual más correlacionada con el semestre de término fue en ambas carreras
\texttt{s8\_reg} ($r = -0.818$ en Mat, $r = -0.785$ en Fís), anticipando el hallazgo dominante
de los modelos predictivos posteriores.

\subsection{Análisis de Supervivencia (Capítulo 6)}

La prueba Log-Rank confirmó una diferencia estadísticamente significativa entre las curvas de
supervivencia de Matemáticas y Física ($\chi^2 = 5.85$, $p = 0.016$): la mediana de egreso en
Física es de 17.0 semestres $[16.0, 18.0]$, mientras que en Matemáticas más del 50\,\% de la
cohorte no alcanza la titulación en el horizonte de observación de 20 semestres.

\noindent
La segmentación por perfil académico al cuarto semestre demostró que la paradoja del
perfeccionismo no es marginal: la Log-Rank entre Pragmáticos y Perfeccionistas arrojó
$p < 0.0001$, y el modelo de Cox asignó al perfil de regularidad alta un Hazard Ratio de
$12.70$ — un estudiante con alta regularidad tiene $12.70$ veces más probabilidad instantánea
de egresar que uno con baja regularidad, en igualdad de condiciones de promedio.

\subsection{Modelos Predictivos (Capítulo 7)}

\subsubsection{Enfoque Multiclase --- Regresión Logística Multinomial}

El sistema de alerta temprana (Modelo~A, S1--S4) clasificó correctamente el estado académico
final del $58.4\,\%$ de los estudiantes. El modelo de
trayectoria completa (Modelo~B, S1--S8) incrementó significativamente la exactitud global a
$73.6\,\%$.
 La prueba de McNemar confirmó que esta diferencia de
$+15.2$ puntos es estadísticamente significativa ($\chi^2 = 639.26$,
$p = 4.83 \times 10^{-141}$), no un artefacto del muestreo.

\subsubsection{Enfoque Truncado — Regresión Regularizada}

Sobre la submuestra libre de censura ($N = 3{,}157$, generaciones $\leq 2016$), Ridge y Lasso
alcanzaron $R^2 \approx 0.755$, prediciendo el semestre exacto de egreso con un Error Absoluto
Medio de $0.957$ semestres. La prueba de Diebold-Mariano no encontró diferencia significativa
entre ambos ($DM = 0.026$, $p = 0.979$), por lo que se prefiere Lasso por su mayor parsimonia
(31 vs.\ 40 predictores activos). El único predictor con $|t| > 10$ fue \texttt{s8\_reg}
($\beta = -2.445$, $t = -13.64$), consolidando la conclusión central.

% ──────────────────────────────────────────────────────────────────────────────
\section{Implicaciones Institucionales}
% ──────────────────────────────────────────────────────────────────────────────

Los hallazgos de esta investigación tienen implicaciones directas para la gestión académica
de la Facultad de Ciencias:

\begin{enumerate}
    \item \textbf{Reorientar la métrica de riesgo.} Los sistemas de alerta estudiantil actuales
          suelen basarse en el promedio de calificaciones. Esta investigación demuestra que la
          \emph{regularidad semestral} —es decir, el cumplimiento de créditos en tiempo y
          forma— es un predictor $12.7\times$ más potente del riesgo de rezago. Los paneles
          de seguimiento deben incorporar un indicador de regularidad acumulada como señal
          de alerta primaria.

    \item \textbf{Intervención temprana en el segundo año.} La divergencia de trayectorias
          comenzó a ser estadísticamente significativa desde el segundo y tercer semestre.
          Un sistema de alerta basado en los primeros cuatro semestres ya logra una precisión
          del $65.5\,\%$ en la clasificación del destino académico final. Esto justifica
          programas de tutoría y apoyo focalizados en los estudiantes con baja regularidad
          temprana, independientemente de su promedio.

    \item \textbf{Redefinir el éxito académico.} La estrategia cultural de priorizar las
          calificaciones sobre la continuidad —frecuentemente incentivada por becas y sistemas
          de evaluación que premian el promedio— puede ser contraproducente para el egreso
          oportuno. La Facultad podría considerar mecanismos que incentiven la inscripción
          continua al mismo nivel que los estímulos por promedio.

    \item \textbf{Capacidad de predicción cuantitativa.} El modelo Lasso del Enfoque Truncado
          puede predecir el semestre de egreso de un estudiante con un error menor a un semestre
          ($\text{MAE} = 0.957$). Integrado en un sistema de información escolar, este modelo
          permitiría generar proyecciones personalizadas de egreso y optimizar la asignación
          de recursos de apoyo académico.
\end{enumerate}

% ──────────────────────────────────────────────────────────────────────────────
\section{Limitaciones y Líneas de Investigación Futura}
% ──────────────────────────────────────────────────────────────────────────────

\subsection{Limitaciones del Estudio}

\begin{itemize}
    \item \textbf{Alcance institucional.} Los datos corresponden exclusivamente a las
          licenciaturas en Matemáticas y Física. La generalización de los hallazgos a otras
          disciplinas o instituciones requiere validación en nuevos contextos.

    \item \textbf{Variables no observadas.} El Vector Excalibur no captura factores
          socioeconómicos, personales o de salud que también influyen en el rezago. La
          inercia académica es el mecanismo \emph{estadísticamente dominante} en los datos
          disponibles, pero no necesariamente el único causal.

    \item \textbf{Clase ``Activo/Aún cursando''.} Ambos modelos de clasificación tienen
          dificultad para identificar a los estudiantes activos (recall $< 5\,\%$). Esto es
          inherente al problema: su trayectoria no ha concluido y se solapa con la de quienes
          se titularán o abandonarán. Datos longitudinales actualizados semestralmente
          resolverían esta limitación.

    \item \textbf{Supuesto de censura no informativa.} El análisis de supervivencia asume
          que el mecanismo de censura es independiente del tiempo de egreso. La ausencia de
          datos sobre las causas reales de abandono impide validar formalmente este supuesto.
\end{itemize}

\subsection{Líneas de Investigación Futura}

\begin{itemize}
    \item \textbf{Modelos de tiempo discreto.} Explorar modelos de supervivencia paramétricos
          con distribuciones de Weibull o Log-Normal para mejorar la predicción del tiempo
          de egreso más allá de las generaciones $\leq 2016$.

    \item \textbf{Aprendizaje profundo sobre secuencias.} Dado el carácter secuencial de la
          trayectoria semestral, modelos de redes recurrentes (LSTM, Transformer) o de
          proceso de punto temporal podrían capturar dependencias temporales no lineales que
          la regresión regularizada no modela explícitamente.

    \item \textbf{Extensión a otras licenciaturas.} Replicar el análisis en Biología, Física
          Biomédica u otras disciplinas de la Facultad de Ciencias permitiría identificar
          si la paradoja del perfeccionismo es universal o específica de las carreras
          de Matemáticas y Física.

    \item \textbf{Actualización en tiempo real.} Diseñar un pipeline de ingesta semestral
          que actualice las predicciones del modelo Lasso conforme se dispone de nuevos
          datos, habilitando un sistema de alerta dinámico y prospectivo.
\end{itemize}

% ──────────────────────────────────────────────────────────────────────────────
\section{Reflexión Final}
% ──────────────────────────────────────────────────────────────────────────────

Esta investigación ofrece una respuesta cuantitativa a una pregunta institucional urgente:
¿\emph{quién llega y quién no?} Los datos demuestran que el egreso oportuno no es función
del talento medido en calificaciones, sino de la \textbf{persistencia medida en regularidad}.
El estudiante que se inscribe semestre a semestre, aunque no brille en promedio, tiene
sistemáticamente mejores probabilidades de concluir su carrera que aquel que sacrifica
materias en busca de la excelencia en la nota.

Esta conclusión no es trivial. Sugiere que el diseño institucional —planes de estudio,
políticas de beca, sistemas de seguimiento— debe migrar de una lógica de \emph{control de
calidad ex post} (intervenir cuando el promedio cae) a una lógica de \emph{monitoreo de
continuidad en tiempo real} (intervenir cuando la regularidad disminuye). En ese sentido,
los modelos predictivos desarrollados en esta tesis no son solo un ejercicio académico:
son prototipos de herramientas de gestión escolar basadas en evidencia estadística rigurosa.
